\section{问题分析}

三问的难点并不是分别找到三个孤立算法，而是控制同一调度决策在逻辑生命周期、物理内存与时间流水三个层次上的连锁影响。

\subsection{问题一：调度顺序决定生命周期重叠}

同一个 DAG 通常存在大量合法拓扑序。某个普通计算节点本身不会直接改变 L1+UB 驻留量，却可能解锁一个大 buffer 的 FREE；某个新的 ALLOC 虽然当前可执行，却可能过早打开生命周期并抬高后续峰值。因此问题一不能只看节点当前的内存增量，而应在保持拓扑与 L0 合法性的同时，比较不同 ready-set 决策对未来释放机会的影响。

完整枚举拓扑序对最大实例不可行。故本文选择“多个确定性启发式 + 有界前瞻 + 独立 evaluator”的策略：启发式负责高效地产生少量结构互补的合法候选，统一评价器负责最终择优。

\subsection{问题二：连续地址使总容量约束升级为碎片约束}

问题二并非简单检查
\[
\sum_{b\in\mathrm{resident}}z_b\le C_t.
\]
即使总空闲容量足够，也可能没有足够长的连续空闲区间容纳新 buffer。于是需要同时处理动态空闲块合并、地址复用、SPILL victim 选择以及 reload 后的新地址分配。

进一步地，不同 buffer 的官方 SPILL traffic 代价并不相同；因此“换出最大的 buffer”“换出最久不用的 buffer”或“最小化 SPILL 次数”都可能偏离题面目标。本文把问题转化为：对每次无法分配的请求直接寻找一个可清空的连续目标窗口，并以清空窗口的官方字节级 traffic 为第一优化目标。

\subsection{问题三：更紧凑的内存布局不一定更快}

物理地址复用会引入额外的执行依赖。同一地址先后服务两个 buffer 时，后者必须等待前者对应物理 residency 结束，因此 Q2 中有利于减少 SPILL 的紧凑布局可能压缩 Q3 的并行空间。另一方面，MTE1/2/3、CUBE、VECTOR、FIXP 等 Pipe 又允许跨流水并行，使全局 schedule 中相邻的节点并不等价于时间轴上的串行执行。

因此问题三需要显式构造包含 DAG、SPILL、地址复用与 Pipe serialization 的 augmented precedence graph，并在这一时间模型上寻找关键路径瓶颈。为了不混淆时间收益和搬运成本，本文首先固定 Q2 的 traffic/SPILL 决策优化 official cycles，再单独讨论允许少量额外 traffic 时的 Pareto 权衡。

\subsection{总体策略}

基于以上分析，本文采用逐层递进的求解路线：
\[
\boxed{\text{合法拓扑序}}
\longrightarrow
\boxed{\text{连续地址 + SPILL}}
\longrightarrow
\boxed{\text{Pipe 时间优化}}.
\]
每一层都保留上一层的硬约束，并通过独立重放器检查输出。这样既可以在大规模实例上使用启发式搜索，又能把最终定量结论与候选生成算法本身解耦。